\subsection{应用：II. 测度的函数}

设 \(\mu_1, \mu_2, \ldots, \mu_n\) 是 \(T\) 上的实测度，并设 \(u(x_1, \ldots, x_n)\) 是定义在 \(\mathbf{R}^n\) 上的一个有限数值函数，且是正齐次的，也就是说（Ch. I, §1, No.~1），满足

\[
u (\alpha x _ {1}, \ldots , \alpha x _ {n}) = \alpha u (x _ {1}, \ldots , x _ {n})\tag{12}
\]

对每个标量 \(\alpha \geqslant 0\) 。存在 T 上的正测度 \(\lambda\) ，使得 \(|\mu_i| \leqslant \lambda\) 对 \(1 \leqslant i \leqslant n\) 成立（例如和 \(\sum_{i=1}^{n} |\mu_i|\) ）。设 \(\lambda\) 和 \(\lambda'\) 是 T 上的两个这样的测度。可以写 \(\mu_i = f_i \cdot \lambda = f_i' \cdot \lambda'\) ，其中 \(f_i\) （相应地，\(f_i'\) ）对测度 \(\lambda\) （相应地，\(\lambda'\) ）是可测的且本性有界的（No.~5, 定理 2）。我们将建立下列结果：为使数值函数 \(u(f_1, \ldots, f_n)\) 对 \(\lambda\) 局部可积，必须且只须函数 \(u(f_1', \ldots, f_n')\) 对 \(\lambda'\) 局部可积，此时

\[
u (f _ {1}, \dots , f _ {n}) \cdot \lambda = u (f _ {1} ^ {\prime}, \dots , f _ {n} ^ {\prime}) \cdot \lambda^ {\prime}.
\]

由于 \(|\mu_i| \leqslant \inf(\lambda, \lambda')\) ，我们可以限于 \(\lambda \leqslant \lambda'\) 的情形。此时 \(\lambda = g \cdot \lambda'\) ，其中 \(g\) 是一个 \(\lambda'\) -可测函数，满足 \(0 \leqslant g \leqslant 1\) （No.~5, 定理 2）；由此（No.~4, 命题 8）得

\[
\mu_ {i} = f _ {i} \cdot (g \cdot \lambda^ {\prime}) = (f _ {i} g) \cdot \lambda^ {\prime};
\]

由此（No.~3, 命题 3 的推论 2）推出，\(f_{i}g\) 等于 \(f_{i}'\) （关于 \(\lambda'\) 局部几乎处处）。于是，由 (12)，

\[
u (f _ {1} ^ {\prime}, \dots , f _ {n} ^ {\prime}) = u (f _ {1} g, \dots , f _ {n} g) = u (f _ {1}, \dots , f _ {n}) g
\]

关于 \(\lambda'\) 局部几乎处处成立。为使 \(u(f_1', \ldots, f_n')\) 是局部 \(\lambda'\) -可积的，因此必须且只须 \(u(f_1, \ldots, f_n)g\) 对 \(\lambda'\) 局部可积，从而（No.~4, 命题 8）只须 \(u(f_1, \ldots, f_n)\) 对 \(\lambda\) 局部可积；此时（No.~4, 命题 8）

\[
u \left(f _ {1} ^ {\prime}, \dots , f _ {n} ^ {\prime}\right) \cdot \lambda^ {\prime} = \left(u \left(f _ {1}, \dots , f _ {n}\right) g\right) \cdot \lambda^ {\prime} = u \left(f _ {1}, \dots , f _ {n}\right) \cdot \lambda .
\]

于是，测度 \(u(f_{1},\ldots,f_{n})\cdot\lambda\) 只依赖于测度 \(\mu_{1},\ldots,\mu_{n}\) 和函数 u；它也记作 \(u(\mu_{1},\ldots,\mu_{n})\) 。因此，只要 u 是这样一个正齐次函数：对一个正测度 \(\lambda\) （它是所有 \(|\mu_{i}|\) 的上界）而言，\(u(f_{1},\ldots,f_{n})\) 是局部 \(\lambda\) -可积的（其中 \(f_{i}\) 表示 \(\mu_{i}\) 关于 \(\lambda\) 的密度），这个测度就有定义。注意，当 u 是正齐次且连续的函数时，这一条件是满足的：因为此时有

\[
\left| u \left(x _ {1}, \dots , x _ {n}\right) \right| \leqslant a \left(\left| x _ {1} \right| + \left| x _ {2} \right| + \dots + \left| x _ {n} \right|\right)
\]

（u 在 \((0,\ldots,0)\) 的一个充分小的邻域中有界），而且由于 \(u(f_{1},\ldots,f_{n})\) 是 \(\lambda\) -可测的（Ch. IV, §5, No.~3, 定理 1），根据可积性判别法（Ch. IV, §5, No.~6, 定理 5），它是局部 \(\lambda\) -可积的。

设 \(u_1, \ldots, u_p\) 是定义在 \(\mathbf{R}^n\) 上的正齐次数值函数，使得这 \(p\) 个函数 \(g_k = u_k(f_1, \ldots, f_n)\) （\(1 \leqslant k \leqslant p\) ）都是局部 \(\lambda\) -可积的。设 \(v\) 是定义在 \(\mathbf{R}^p\) 上的一个正齐次数值函数，使得 \(v(g_1, \ldots, g_p)\) 是局部 \(\lambda\) -可积的。置

\[
w (x _ {1}, \ldots , x _ {n}) = v \big (u _ {1} (x _ {1}, \ldots , x _ {n}), \ldots , u _ {p} (x _ {1}, \ldots , x _ {n}) \big).
\]

于是，函数 w 是正齐次的，\(w(f_{1}, \ldots, f_{n})\) 是局部 \(\lambda\) -可积的，且依定义，

\[
w \left(\mu_ {1}, \dots , \mu_ {n}\right) = v \left(u _ {1} \left(\mu_ {1}, \dots , \mu_ {n}\right), \dots , u _ {p} \left(\mu_ {1}, \dots , \mu_ {n}\right)\right).
\]

在函数 \(x^{+}\) ，\(x^{-}\) ，\(|x|\) ，\(x + y\) ，\(\inf(x, y)\) ，\(\sup(x, y)\) 的特殊情形，刚才所述程序所定义的测度分别与已记作 \(\mu^{+}\) ，\(\mu^{-}\) ，\(|\mu|\) ，\(\mu + \nu\) ，\(\inf(\mu, \nu)\) ，\(\sup(\mu, \nu)\) 的那些测度重合；这由 No.~2 的命题 2 的推论立即得出。若 \(\mu\) 和 \(\nu\) 是两个实测度，且 \(\theta = \mu + i\nu\) ，则 \(|\theta| = \sqrt{\mu^{2} + \nu^{2}}\) ；因为，设 \(\lambda\) 是一个测度 \(\geqslant 0\) ，它是 \(|\mu|\) 和 \(|\nu|\) 的一个上界，并设 f, g 是局部 \(\lambda\) -可积函数，使得 \(\mu = f \cdot \lambda\) ，\(\nu = g \cdot \lambda\) ；则

\[
\sqrt {\mu^ {2} + \nu^ {2}} = \sqrt {f ^ {2} + g ^ {2}} \cdot \lambda ,
\]

\(\theta = (f + ig) \cdot \lambda\) ，因此（No.~2, 命题 2）\(|\theta| = \sqrt{f^2 + g^2} \cdot \lambda\) 。

这个方法可以应用于正齐次函数 \((x_1^2 +\ldots +x_n^2)^{1 / 2}\) ，以定义 \(\mathbf{R}^n\) 中一条曲线的长度

